\subsection{Простые схемы}
%%%%%%%%%%%%%%%%%%%%%%%%%%%%%%%%%%%%%%%%%%%%%%%%%%%%
%%%%%%%%%%%%%%%%%%%%%%%%%%%%%%%%%%%%%%%%%%%%%%%%%%%%
 \z[{\begin{wrapfigure}[3]{r}{0.25\textwidth}%
		\vspace{-30pt}%
		\begin{flushright}%
 			\includegraphics[width=0.25\textwidth]{3.4.1}%
		\end{flushright}%
	\end{wrapfigure}%
	}]%
	{Какой ток проходит через каждое сопротивление, если $R_1 = R_2 = R_3 = R_4 = 1\,\text{Ом}$, $R_5 = 3\,\text{Ом}$, напряжение $U = 12\,\text{В}$.
	\vspace{1em}
	}
%
{$I_1 = I_2 = I_3 = I_4 = \frac{U}{2R} = 6\,\text{А};\quad I_5 = 0$}

%%%%%%%%%%%%%%%%%%%%%%%%%%%%%%%%%%%%%%%%%%%%%%%%%%%%
 \z[{\begin{wrapfigure}[6]{r}{0.25\textwidth}%
		\vspace{-15pt}%
		\begin{flushright}%
 			\includegraphics[width=0.25\textwidth]{3.4.2}%
		\end{flushright}%
	\end{wrapfigure}%
	}]%
	{Для схемы, изображенной на рисунке, подобрать такое сопротивление $R$, чтобы ток, текущий через это сопротивление при замкнутом ключе $K_1$ и разомкнутом ключе $K_2$, был в $3$ раза больше тока, текущего через это сопротивление при разомкнутом ключе $K_1$ и замкнутом ключе $K_2$. Сопротивлением источника пренебречь.
	\vspace{1em}
	}
%
{$R = r$}
%%%%%%%%%%%%%%%%%%%%%%%%%%%%%%%%%%%%%%%%%%%%%%%%%%%%
 \z[{\begin{wrapfigure}[4]{r}{0.35\textwidth}%
		\vspace{-10pt}%
		\begin{flushright}%
 			\includegraphics[width=0.35\textwidth]{3.4.3}%
		\end{flushright}%
	\end{wrapfigure}%
	}]%
	{Участок цепи, схема которого приведена на рисунке, включает в себя резисторы с сопротивлениями $R$ и $2R$. Амперметр $A_1$ показывает силу тока $I_1 = 0.2\,\text{мА}$. Найдите показания амперметра $A_2$.
	\vspace{1em}
	}
%
{$I_2 = \frac{3}{2}I_1 = 0.3\,\text{мА}$}
%%%%%%%%%%%%%%%%%%%%%%%%%%%%%%%%%%%%%%%%%%%%%%%%%%%%
 \z[{\begin{wrapfigure}[4]{r}{0.25\textwidth}%
		\vspace{-25pt}%
		\begin{flushright}%
 			\includegraphics[width=0.25\textwidth]{3.4.4}%
		\end{flushright}%
	\end{wrapfigure}%
	}]%
	{Идеальный вольтметр включён в цепь, схема которой изображена на рисунке. Цепь состоит из четырёх одинаковых резисторов сопротивлением $R$ и батареи с ЭДС $\emf = 9\,\text{В}$ и нулевым внутренним сопротивлением. Найдите показания вольтметра.
	\vspace{1em}
	}
%
{$V = \frac{\emf}{5}$}
%%%%%%%%%%%%%%%%%%%%%%%%%%%%%%%%%%%%%%%%%%%%%%%%%%%%
 \z[{\begin{wrapfigure}[4]{r}{0.25\textwidth}%
		\vspace{-35pt}%
		\begin{flushright}%
 			\includegraphics[width=0.25\textwidth]{3.4.5}%
		\end{flushright}%
	\end{wrapfigure}%
	}]%
	{Идеальный амперметр в цепи, схема которой изображена на рисунке, показывает силу тока $I = 9\,\text{мА}$. Определите сопротивление резистора $R$, если напряжение идеального источника $U = 6\,\text{В}$.
	\vspace{1em}
	}
%
{$R = \frac{9}{10} \cdot \frac{U}{I} = 600\,\text{Ом}$}

%%%%%%%%%%%%%%%%%%%%%%%%%%%%%%%%%%%%%%%%%%%%%%%%%%%%
 \z[{\begin{wrapfigure}[6]{r}{0.25\textwidth}%
		\vspace{-10pt}%
		\begin{flushright}%
 			\includegraphics[width=0.25\textwidth]{3.4.6}%
		\end{flushright}%
	\end{wrapfigure}%
	}]%
	{На рисунке изображена схема электрической цепи, состоящей из источника постоянного напряжения $U_0$, резисторов с одинаковым сопротивлением $R$, идеального вольтметра и идеального амперметра. Показания вольтметра $U_V = 3\,\text{В}$, амперметра — $I_A = 24\,\text{мА}$. Определите напряжение источника $U_0$ и сопротивление $R$ резисторов. Сопротивление источника считать равным нулю.
	\vspace{1em}
	}
%
{$U_0 = \frac{5}{3}U_V = 5\,\text{В};\quad R = \frac{6}{5} \cdot \frac{U_0}{I_A} = 250\,\text{Ом}$}

%%%%%%%%%%%%%%%%%%%%%%%%%%%%%%%%%%%%%%%%%%%%%%%%%%%%
 \z[{\begin{wrapfigure}[4]{r}{0.25\textwidth}%
		\vspace{+15pt}%
		\begin{flushright}%
 			\includegraphics[width=0.25\textwidth]{3.4.7}%
		\end{flushright}%
	\end{wrapfigure}%
	}]%
	{Электрическая цепь, схема которой изображена на рисунке, подключена к батарейке. Вольтметры $V_1$ и $V_2$ показывают напряжения $U_1 = 1\,\text{В}$ и $U_2 = 0.1\,\text{В}$, а амперметр $A$ показывает силу тока $I = 1\,\text{мА}$. Найдите сопротивления приборов. Вольтметры считайте одинаковыми.}
%
{$R_A = \frac{U_2}{I} = 100\,\text{Ом};\quad R_V = \frac{U_1 - U_2}{I} = 900\,\text{Ом}$}


